% Seleciona o idioma do documento (conforme pacotes do babel)
%\selectlanguage{english}
\selectlanguage{brazil}

% Retira espaço extra obsoleto entre as frases.
\frenchspacing 

\newpage

% ==============================================
% ELEMENTOS PRÉ-TEXTUAIS
% ==============================================
\pretextual

% ----------------------------------------------
% Capa
% ----------------------------------------------
%\imprimircapa
% Capa personalizada sem o uso de \imprimircapa
\begin{capa} 
    \begin{center}
        \begin{minipage}{1\textwidth} 
            \large\centering\makebox[\textwidth]{\includegraphics[scale=1]{imagens/Logo UNIP.png}}
        \end{minipage}
    \end{center}
    \begin{center}
        \LARGE\textbf{UNIP – UNIVERSIDADE PAULISTA\\} 
        \LARGE Curso de Ciência da Computação\\
        \vfill
        \ABNTEXchapterfont\Large\textbf{\MakeUppercase{\imprimirtitulo}}
        \\\small{AS TÉCNICAS CRIPTOGRÁFICAS, CONCEITOS, USOS E APLICAÇÕES\\UTILIZANDO A TÉCNICA TRIPLE DES} 
        \vfill
        \normalsize
        Daniel Chrispim - F263614\\
        Daniel Lima - N621184\\
        Gabriel De Paula Seki - N6853H6\\
        Maisa Serpa Castro Moreira - F17EFC7\\
        Matheus Lopes - N620ED0\\
        Vinícius Amorim - N641JC
        \vfill
        \imprimirlocal, \imprimirdata 
    \end{center}
\end{capa}

% ----------------------------------------------
% Folha de rosto
% ----------------------------------------------
% folha de rosto personalizada sem uso de \imprimirfolhaderosto
\makeatletter
\renewcommand{\folhaderostocontent}{
\begin{center}
    \begin{center}
        \begin{minipage}{1\textwidth} 
            \large\centering\makebox[\textwidth]{\includegraphics[scale=1]{imagens/Logo UNIP.png}}
        \end{minipage}
    \end{center}
    
  \vspace*{\fill}%\vspace*{\fill}
  \begin{center}
  \ABNTEXchapterfont\Large\textbf{\MakeUppercase{\imprimirtitulo}}
  \\\small{AS TÉCNICAS CRIPTOGRÁFICAS, CONCEITOS, USOS E APLICAÇÕES\\UTILIZANDO A TÉCNICA TRIPLE DES}
  \end{center}
  \vspace*{\fill}
  
    \hspace{.45\textwidth}
    \begin{minipage}{.5\textwidth}
    \SingleSpacing
    {Atividades Práticas Supervisionadas do 2\degree\ Semestre do Curso de Ciência da Computação da \textbf{Universidade Paulista UNIP.}}
    \end{minipage}%
    \vspace*{\fill}


  \hspace{.45\textwidth}
  \begin{minipage}{.5\textwidth}
	{\textbf{Coordenador:} Prof. Fernando A. Gotti}%
  \end{minipage}%
 
  
  \hspace{.45\textwidth}
  \begin{minipage}{.5\textwidth}
	{\textbf{Prof. Responsável:} Luiz Gustavo Miranda Pinto}%
  \end{minipage}%

  
  \vspace*{\fill}
  %{\abntex@ifnotempty{\imprimirinstituicao}{\imprimirinstituicao\vspace*{\fill}}}

  \imprimirlocal, \imprimirdata
\end{center}
}
\makeatother

% Folha de rosto (o * indica que haverá a ficha bibliográfica)

\imprimirfolhaderosto

% \imprimirfolhaderosto*

% ----------------------------------------------
% Inserir a ficha bibliografica catalográfica
% ----------------------------------------------
% Isto é um exemplo de Ficha Catalográfica, ou ``Dados internacionais de catalogação-na-publicação''. Você pode utilizar este modelo como referência. Porem, provavelmente a biblioteca da sua universidade lhe fornecerá um PDF com a ficha catalográfica definitiva após a defesa do trabalho. Quando estiver com o documento, salve-o como PDF no diretório do seu projeto e substitua todo o conteúdo de implementação deste arquivo pelo comando abaixo:

% \begin{fichacatalografica}
%     \includepdf{fig_ficha_catalografica.pdf}
% \end{fichacatalografica}
% \begin{fichacatalografica}
% 	\sffamily
% 	\vspace*{\fill}					% Posição vertical
%   	\begin{center}					% Minipage Centralizado
% 	\fbox{
%     \begin{minipage}[t]{1,5cm} 
%     \vspace{0.5cm} 
%     D383m %Algum número que o bibliotecario ira gerar
%     \end{minipage}
    
%     \begin{minipage}[t]{11cm}	% Largura
% 	\small
%     \vspace{0.5cm}
% 	%\imprimirautor		% ATENCAO - SUBSTITUIR POR %Sobrenome, Nome do autor
%     AutorSobrenome, AutorNome
	
% 	\hspace{0.5cm} 
%     \imprimirtitulo  / \imprimirautor. -- \imprimirdata.
	
% 	\hspace{0.5cm} 
%     \pageref{LastPage} p. : il. (algumas color.) ; 30 cm.\\
	
%     \hspace{0.5cm}
% 	\imprimirtipotrabalho~--~Universidade do Vale do Rio dos Sinos, Programa de Pós-Graduação em Engenharia Elétrica, São Leopoldo, RS, \imprimirdata. \hfill 
    
%     \hspace{0.5cm}
%     \imprimirorientadorRotulo~\imprimirorientador ;
%     \imprimircoorientadorRotulo~\imprimircoorientador\\
	
% 	\hspace{0.5cm}
% 		1. Palavra-chave1
% 		2. Palavra-chave1
% 		3. Palavra-chave1
%         I. Título \\
% 		%I. \imprimirtitulo \\
%         %II. \imprimirorientador \\
    
% 	\hspace{8.75cm} CDU 621.3 %algum outro numero
	
%     \end{minipage}}
    
%     \hspace{0.5cm}
%     Dados Internacionais de Catalogação na Publicação (CIP) \\  	
%     (Bibliotecário: Nome Sobrenome – CRB 10/1298)
	
%     \end{center}
% \end{fichacatalografica}

% ----------------------------------------------
% Inserir errata
% ----------------------------------------------
% \begin{errata}
% Elemento opcional da \citeonline[4.2.1.2]{NBR14724:2011}. Exemplo:

% \vspace{\onelineskip}

% FERRIGNO, C. R. A. \textbf{Tratamento de neoplasias ósseas apendiculares com reimplantação de enxerto ósseo autólogo autoclavado associado ao plasma rico em plaquetas}: estudo crítico na cirurgia de preservação de membro em cães. 2011. 128 f. Tese (Livre-Docência) - Faculdade de Medicina Veterinária e Zootecnia, Universidade de São Paulo, São Paulo, 2011.

% \begin{table}[htb]
% \center
% \footnotesize
% \begin{tabular}{|p{1.4cm}|p{1cm}|p{3cm}|p{3cm}|}
%   \hline
%    \textbf{Folha} & \textbf{Linha}  & \textbf{Onde se lê}  & \textbf{Leia-se}  \\
%     \hline
%     1 & 10 & auto-conclavo & autoconclavo\\
%    \hline
% \end{tabular}
% \end{table}

% \end{errata}

% ----------------------------------------------
% Inserir folha de aprovação
% ----------------------------------------------
% Isto é um exemplo de Folha de aprovação, elemento obrigatório da NBR 14724/2011 (seção 4.2.1.3). Você pode utilizar este modelo até a aprovação do trabalho. Após isso, substitua todo o conteúdo deste arquivo por uma imagem da página assinada pela banca com o comando abaixo:
%
% \includepdf{folhadeaprovacao_final.pdf}
%
% \begin{folhadeaprovacao}
% % Gerada conforme instrucoes do PPGEE
%   \begin{center}
%   {\ABNTEXchapterfont\large\imprimirautor}

%   \vspace*{\fill}\vspace*{\fill}
%     \begin{center}
%     	\ABNTEXchapterfont\bfseries\Large\imprimirtitulo
%     \end{center}
%   \vspace*{\fill}

%   \hspace{.45\textwidth}
%     \begin{minipage}{.5\textwidth}
%     	\imprimirpreambulo 
%     \end{minipage}%

%   \vspace*{\fill}
%     \begin{flushleft}
%   	  Aprovado em 11 de Janeiro de 2016. \\
%     \end{flushleft}
%   \vspace*{\fill}
%     BANCA EXAMINADORA:% \imprimirlocal, \today :
%   \end{center}
   
%    %\assinatura{\textbf{\imprimirorientador} \\ Orientador}
%    %\assinatura{\textbf{\imprimircoorientador} \\ Coorientador} 
%    \assinatura{Prof. Dr. Nome Sobrenome -- Unisinos \\ Avaliador}
%    \assinatura{Prof. Dr. Nome Sobrenome -- UFRGS \\ Avaliador Externo}
%    \assinatura{Prof. Dr. Nome Sobrenome -- UERGS \\ Avaliador Externo}
%    %\assinatura{\textbf{Professor} \\ Convidado 4}
     
%     \vspace*{\fill} \vspace*{\fill}
%     \hspace{.4\textwidth}
%     \begin{minipage}{.5\textwidth}
%     	\imprimirorientador~(Orientador) \\
%         \imprimircoorientador~(Coorientador)
%     \end{minipage}%
    
%     \vspace*{\fill}
%     \begin{flushleft}
%     	Visto e permitida a impressão\\
%         \imprimirlocal
%     \end{flushleft}
    
%     \vspace*{\fill}
%     \hspace{.4\textwidth}
%     \begin{minipage}{.5\textwidth}
%     	Prof. Dr. Nome do Coordenador do curso \\
%         Coordenador PPG em Engenharia Elétrica
%     \end{minipage}%
% %    \begin{center}
% %     \vspace*{0.5cm}
% %       {\large\imprimirlocal}
% %       \par
% %       {\large\imprimirdata}
% %       \vspace*{1cm}
% %   \end{center}
  
% \end{folhadeaprovacao}

% ----------------------------------------------
% Dedicatória
% ----------------------------------------------
% \begin{dedicatoria}
%    \vspace*{\fill}
%    \centering
%    \noindent
%    \textit{ Este trabalho é dedicado às crianças adultas que,\\
%    quando pequenas, sonharam em se tornar cientistas.} 
%    \vspace*{\fill}
% \end{dedicatoria}

% ----------------------------------------------
% Agradecimentos
% ----------------------------------------------
% \begin{agradecimentos}
%     \lipsum[1]
% \end{agradecimentos}

% ----------------------------------------------
% Epígrafe
% ----------------------------------------------
% Importante: O autor da epígrafe deve constar na lista de referências
% \begin{epigrafe}
%     \vspace*{\fill}
% 	\begin{flushright}
% 		\textit{``Os que se encantam com a prática sem a ciência \\
%         são como os timoneiros que entram no navio sem timão nem bússola,\\
%         nunca tendo certeza do seu destino."\\
%         (Leonardo da Vinci)}
% 	\end{flushright}
% \end{epigrafe}

% ||||||||||||||||||||||||||||||||||||||||||||||
% RESUMOS
% ||||||||||||||||||||||||||||||||||||||||||||||

% ----------------------------------------------
% Resumo em português
% ----------------------------------------------
% Importante: De acordo com a NBR6024 as palavras-chaves devem ser separadas entre si por ponto e devem ter somente a primeira palavra escrita com letra maiúscula
% \setlength{\absparsep}{18pt} % ajusta o espaçamento dos parágrafos do resumo
% \begin{resumo}
    
% 	\vspace{\onelineskip}
 
% 	\noindent 
% 	\textbf{Palavras-chaves}: Sistemas microeletromecânicos (MEMS). Isoladores e isolamentos elétricos. Fontes Chaveadas. Fonte chaveada isolada. Isolação elétrica. Fonte de alimentação de baixo custo.
% \end{resumo}

% ----------------------------------------------
% Resumo em inglês
% ----------------------------------------------
% Importante: De acordo com a NBR6024 as palavras-chaves devem ser separadas entre si por ponto e devem ter somente a primeira palavra escrita com letra maiúscula
% \begin{resumo}[Abstract]
% \begin{otherlanguage*}{english}
% 	\lipsum[8]
    
% 	\vspace{\onelineskip}
 
% 	\noindent 
% 	\textbf{Key-words}: Microelectromechanical systems (MEMS). Electronic insulation. Switching power supply. Isolated power supply. Electric isolated circuitry. Low cost power supply.
% \end{otherlanguage*}
% \end{resumo}

% ----------------------------------------------
% resumo em francês 
% ----------------------------------------------
% Importante: De acordo com a NBR6024 as palavras-chaves devem ser separadas entre si por ponto e devem ter somente a primeira palavra escrita com letra maiúscula
% \begin{resumo}[Résumé]
%  \begin{otherlanguage*}{french}
%     Il s'agit d'un résumé en français.
 
%    \textbf{Mots-clés}: latex. abntex. publication de textes.
%  \end{otherlanguage*}
% \end{resumo}

% ----------------------------------------------
% resumo em espanhol
% ----------------------------------------------
% Importante: De acordo com a NBR6024 as palavras-chaves devem ser separadas entre si por ponto e devem ter somente a primeira palavra escrita com letra maiúscula
% \begin{resumo}[Resumen]
%  \begin{otherlanguage*}{spanish}
%    Este es el resumen en español.
  
%    \textbf{Palabras clave}: latex. abntex. publicación de textos.
%  \end{otherlanguage*}
 
% \end{resumo}

% ----------------------------------------------
% inserir lista de ilustrações
% ----------------------------------------------
% \pdfbookmark[0]{\listfigurename}{lof}
% \listoffigures*
% \cleardoublepage

% Diferentes tipos de listas podem ser criadas por meio de macros do memoir.

% ----------------------------------------------
% inserir lista de tabelas
% ----------------------------------------------
% \pdfbookmark[0]{\listtablename}{lot}
% \listoftables*
% \cleardoublepage

% ----------------------------------------------
% inserir lista de abreviaturas e siglas
% ----------------------------------------------
% Importante: As abreviaturas e siglas devem estar em ordem alfabética
% \begin{siglas}
%   \item[ABNT] Associação Brasileira de Normas Técnicas
%   \item[abnTeX] ABsurdas Normas para TeX
% \end{siglas}

% ----------------------------------------------
% inserir lista de símbolos
% ----------------------------------------------
% Importante: Os símbolos devem estar na ordem de aparecimento no texto.
% \begin{simbolos}
%   \item[$ \Gamma $] Letra grega Gama
%   \item[$ \Lambda $] Lambda
%   \item[$ \zeta $] Letra grega minúscula zeta
%   \item[$ \in $] Pertence
% \end{simbolos}

% ----------------------------------------------
% inserir o sumário
% ----------------------------------------------
\pdfbookmark[0]{\contentsname}{toc}
\tableofcontents*
\cleardoublepage
